% Generated draft: requires mathematical and source review.
\documentclass[11pt]{article}
\usepackage[margin=25mm]{geometry}
\usepackage{fontspec}
\setmainfont{Latin Modern Roman}
\usepackage{amsmath,amssymb,mathtools,unicode-math}
\setmathfont{Latin Modern Math}
\usepackage{enumitem,needspace,xurl,hyperref}
\setlength{\parindent}{0pt}
\setlength{\parskip}{5pt}
\setlength{\emergencystretch}{3em}
\newcommand{\N}{\mathbb{N}}
\newcommand{\Z}{\mathbb{Z}}
\newcommand{\Q}{\mathbb{Q}}
\newcommand{\R}{\mathbb{R}}
\newcommand{\C}{\mathbb{C}}
\newcommand{\F}{\mathbb{F}}
\newcommand{\A}{\mathbb{A}}
\newcommand{\PP}{\mathbb{P}}
\newcommand{\E}{\mathbb{E}}
\newcommand{\eps}{\varepsilon}
\newcommand{\dd}{\,\mathrm{d}}
\begin{document}
\section*{UnsolvedMath ID: 1100810}
\textbf{Source problem code:} catalogue entry
\subsection{Definitions and mathematical statement}
Let \(C\) be a nontrivial finite group, and let \(H=F/N\) be a finitely presented group containing a subgroup isomorphic to \(C^n\) for every integer \(n\ge1\), where \(C^n\) denotes the direct product of \(n\) copies of \(C\). In the finite-presentation interpretation, \(F\) is a fixed finite-rank free group and \(N\trianglelefteq F\) is its presentation kernel. Choose embeddings \(\iota_n:C^n\hookrightarrow H\), and amalgamate two copies of \(H\) using the corresponding copies of \(\iota_n\):
\[
G_n=H*_{C^n}H,\qquad
N_n=\ker(F*F\longrightarrow G_n),\qquad
G_n\cong(F*F)/N_n.
\]
For \(K\trianglelefteq L\), let \(d_L(K)\) be the minimum cardinality of a subset of \(K\) whose normal closure in \(L\) equals \(K\). The question asks whether, for every such setup,
\[
\lim_{n\to\infty}d_{F*F}(N_n)=\infty,
\quad\text{that is,}\quad
\forall R\in\N\ \exists n_R\ \forall n\ge n_R:\ d_{F*F}(N_n)\ge R.
\]
The embeddings and their possible compatibility as \(n\) varies are not specified in the excerpt. The displayed construction uses the same chosen embedding in the two copies of \(H\); an interpretation using different embeddings should be identified separately.

A finite group \(C\) is perfect when \(C=[C,C]\). The supplied partial-progress summary claims an affirmative answer for every nonperfect finite \(C\), via linear growth in the number of generators of \(H_2(G_n;\Z)\) and a homological relator lower bound. It attributes this to a new argument ``above,'' which is not included in the record. It reports the perfect case open, identifies that case as the one relevant to the relation-gap problem and the source construction, and says the relation-gap problem remained open in literature checked through December 2024. Its further assertions that known computable lower bounds stay bounded in the perfect case and that this case matches the difficulty of Bridson--Tweedale's ``putative relation gap'' are supplied assessments, not verified equivalences here.
\subsection{Short English statement}
For the doubles of a fixed finitely presented group amalgamated along larger direct powers of a finite group, must the normal number of defining relations grow without bound? The supplied summary claims the nonperfect case and leaves the perfect case open, but its claimed new argument is absent.
\subsection{Sources}
\begin{itemize}
\item Local UnsolvedMath registry, permanent ID \(1100810\), exact source code \texttt{catalogue entry}, ``Questions in Geometric Group Theory --- Q 8.10.''
\item Bestvina, \emph{Questions in Geometric Group Theory} (2004), Question 8.10, PDF page 18. \url{https://www.math.utah.edu/~bestvina/eprints/questions-updated.pdf}; manual transcription; recorded access 2026-07-29.
\item Supplied partial-progress summary names Bridson--Tweedale and a literature check through December 2024, without complete citations. Its affirmative nonperfect-case argument is referred to as ``above'' but is not supplied; no authoritative reference for that claimed new result is given.
\end{itemize}
\end{document}
